\documentclass[11pt,a4paper]{article}
\usepackage{amsmath,amssymb,amsthm}
\usepackage[margin=27mm]{geometry}
\newtheorem{proposition}{Proposition}
\title{A fresh-contact obstruction to a linear trace incidence}
\author{}\date{17 November 2026}
\begin{document}\maketitle
Write $V=\operatorname{Tr}(W^{-1})$ for the actual detector trace
and $C=V^2+w(2-V)$ for the trace-conic radicand. The relation
$V+w=2$ would make
\[
 C=2\bigl(1+(V-1)^2\bigr),
\]
so multiplying $C$ by that trace class would give a square.
This particular incidence is incompatible with the retained fresh
comparison contacts.
\begin{proposition}
Suppose the actual general or conditional detector tuple retains
a fresh selected comparison contact. Then $V+w\ne2$.
\end{proposition}
\begin{proof}
At that contact $v_p(w)=1$. The successful reciprocal slot is one
modulo $p$. The other is zero: otherwise all detector factors in
both slots would be units, and the two tests $f(H)$ and
$\overline{f(H-P)}$ would force both $H$ and $H-P$ into the formal
group at the same place. Their difference $P$ has finite reduction
at the selected good prime, a contradiction. Thus $V\equiv1\pmod p$,
whereas the proposed incidence would give $V\equiv2\pmod p$.
\end{proof}
The assertion concerns this single incidence on tuples retaining the
specified fresh contact. It does not rule out another trace-class
realization, an isolated rational witness, or a different source
preparation.
\end{document}
